\documentclass[a4paper]{article}

\usepackage{graphicx}
\usepackage{amsmath,amssymb}
\usepackage{minted}
\usepackage{multirow}
\usepackage{float}
\usepackage{todonotes}
\usepackage{forest}
\usepackage{xstring}
\usepackage{xcolor}
\usepackage[most]{tcolorbox}
\usepackage{xparse}
\usepackage{natbib}

\usetikzlibrary{graphs,graphdrawing}
\usegdlibrary{layered}

% for external graphviz
\input{/home/donkarlo/Dropbox/repo/nd_ai_project/src/nd_ai/knoweldge_representation/language/formal/computer/programming/tex/latex/code_clips/external_graphviz.tex}

% for tags
\input{/home/donkarlo/Dropbox/repo/nd_ai_project/src/nd_ai/knoweldge_representation/language/formal/computer/programming/tex/latex/parts/control_sequence/kind/semtag/semtag.tex}

% for links
\input{/home/donkarlo/Dropbox/repo/nd_ai_project/src/nd_ai/knoweldge_representation/language/formal/computer/programming/tex/latex/code_clips/seehere.tex}

\title{}
\date{}

\begin{document}
    \maketitle
    
    
    \section{Grammatikthemen im Deutschen {\small (Grammar topics in German)}}
        
        \begin{itemize}
            \item Verben {\small (verbs)}
            \begin{itemize}
                \item Verbformen und Konjugation {\small (verb forms and conjugation)}
                \begin{itemize}
                    \item Konjugation im Präsens {\small (conjugation in present tense)}
                    \item starke und schwache Verben {\small (strong and weak verbs)}
                    \item gemischte Verben {\small (mixed verbs)}
                    \item trennbare und nicht trennbare Verben {\small (separable and inseparable verbs)}
                    \item reflexive Verben {\small (reflexive verbs)}
                    \item Modalverben {\small (modal verbs)}
                \end{itemize}
                
                \item Zeiten {\small (tenses)}
                \begin{itemize}
                    \item Präsens {\small (present)}
                    \item Perfekt {\small (perfect)}
                    \item Präteritum {\small (simple past)}
                    \item Plusquamperfekt {\small (past perfect)}
                    \item Futur I {\small (future I)}
                \end{itemize}
                
                \item Konjunktiv {\small (subjunctive)}
                \begin{itemize}
                    \item Konjunktiv II {\small (subjunctive II)}
                    \begin{itemize}
                        \item hätte, wäre, würde {\small (would have, would be, would)}
                        \item Konjunktiv II der Vergangenheit {\small (subjunctive II of the past)}
                        \item Konjunktiv II mit Modalverben {\small (subjunctive II with modal verbs)}
                        \item irreale Konditionalsätze {\small (unreal conditional clauses)}
                        \item irreale Wunschsätze {\small (unreal wish sentences)}
                        \item Höflichkeit {\small (politeness)}
                        \item vorsichtige Aussagen {\small (tentative statements)}
                    \end{itemize}
                    \item Konjunktiv I {\small (subjunctive I)}
                    \begin{itemize}
                        \item Bildung {\small (formation)}
                        \item Funktionen {\small (functions)}
                    \end{itemize}
                \end{itemize}
                
                \item Verbmuster und Ergänzungen {\small (verb patterns and complements)}
                \begin{itemize}
                    \item Verben mit Akkusativobjekt {\small (verbs with accusative object)}
                    \item Verben mit Dativobjekt {\small (verbs with dative object)}
                    \item Verben mit Dativ- und Akkusativobjekt {\small (verbs with dative and accusative object)}
                    \item Verben mit Präpositionalobjekt {\small (verbs with prepositional object)}
                    \item Verben mit Pronominaladverbien {\small (verbs with pronominal adverbs)}
                    \item Verben mit festen Ergänzungsmustern {\small (verbs with fixed complement patterns)}
                \end{itemize}
                
                \item Verbale Spezialthemen {\small (special verb topics)}
                \begin{itemize}
                    \item Bedeutung und Funktion von \textit{werden} {\small (meaning and function of \textit{werden})}
                    \item \textit{brauchen} und \textit{sich lassen} {\small (verbs \textit{brauchen} and \textit{sich lassen})}
                    \item Nomen-Verb-Verbindungen {\small (noun-verb combinations)}
                    \begin{itemize}
                        \item freie Nomen-Verb-Verbindungen {\small (free noun-verb combinations)}
                        \item Funktionsverbgefüge {\small (light verb constructions)}
                        \item figurative Nomen-Verb-Verbindungen {\small (figurative noun-verb combinations)}
                    \end{itemize}
                \end{itemize}
            \end{itemize}
            
            \item Nomen und Artikel {\small (nouns and articles)}
            \begin{itemize}
                \item Nomen {\small (nouns)}
                \begin{itemize}
                    \item Bedeutung und Form der Nomen {\small (meaning and form of nouns)}
                    \item Genus {\small (gender)}
                    \item Numerus {\small (number)}
                    \item Kasus {\small (case)}
                    \item n-Deklination {\small (n-declension)}
                    \item Komposita {\small (compound nouns)}
                \end{itemize}
                
                \item Artikel {\small (articles)}
                \begin{itemize}
                    \item bestimmte Artikel {\small (definite articles)}
                    \item unbestimmte Artikel {\small (indefinite articles)}
                    \item Nullartikel {\small (zero article)}
                    \item Negation mit \textit{kein} {\small (negation with \textit{kein})}
                \end{itemize}
            \end{itemize}
            
            \item Pronomen und Artikelwörter {\small (pronouns and determiners)}
            \begin{itemize}
                \item Personalpronomen {\small (personal pronouns)}
                \begin{itemize}
                    \item Nominativ {\small (nominative)}
                    \item Akkusativ {\small (accusative)}
                    \item Dativ {\small (dative)}
                \end{itemize}
                
                \item Possessivartikel {\small (possessive determiners)}
                \item Demonstrativartikel und Demonstrativpronomen {\small (demonstrative determiners and pronouns)}
                \item unbestimmte Pronomen und Artikel {\small (indefinite pronouns and articles)}
                \item Pronominaladverbien {\small (pronominal adverbs)}
            \end{itemize}
            
            \item Adjektive und Steigerung {\small (adjectives and comparison)}
            \begin{itemize}
                \item Adjektive {\small (adjectives)}
                \item Adjektivdeklination {\small (adjective declension)}
                \item Komparativ {\small (comparative)}
                \item Superlativ {\small (superlative)}
                \item Vergleichsformen {\small (forms of comparison)}
                \begin{itemize}
                    \item \textit{so \dots wie} {\small (as ... as)}
                    \item \textit{als} {\small (than)}
                \end{itemize}
            \end{itemize}
            
            \item Kasus {\small (cases)}
            \begin{itemize}
                \item Nominativ {\small (nominative)}
                \begin{itemize}
                    \item Subjekt {\small (subject)}
                    \item Prädikativ {\small (predicative)}
                \end{itemize}
                
                \item Akkusativ {\small (accusative)}
                \begin{itemize}
                    \item direktes Objekt {\small (direct object)}
                    \item Akkusativ nach Präpositionen {\small (accusative after prepositions)}
                    \item temporaler Akkusativ {\small (temporal accusative)}
                \end{itemize}
                
                \item Dativ {\small (dative)}
                \begin{itemize}
                    \item indirektes Objekt {\small (indirect object)}
                    \item Dativ nach Präpositionen {\small (dative after prepositions)}
                    \item freier Dativ {\small (free dative)}
                \end{itemize}
                
                \item Genitiv {\small (genitive)}
                \begin{itemize}
                    \item Genitiv als Attribut {\small (genitive as attribute)}
                    \item Genitiv bei Namen {\small (genitive with names)}
                    \item Präpositionen mit Genitiv {\small (prepositions with genitive)}
                \end{itemize}
            \end{itemize}
            
            \item Präpositionen {\small (prepositions)}
            \begin{itemize}
                \item Präpositionen mit Dativ {\small (prepositions with dative)}
                \item Präpositionen mit Akkusativ {\small (prepositions with accusative)}
                \item Präpositionen mit Dativ und Akkusativ {\small (prepositions with dative and accusative)}
                \item Präpositionen mit Genitiv {\small (prepositions with genitive)}
                \item lokale Präpositionen {\small (local prepositions)}
                \item temporale Präpositionen {\small (temporal prepositions)}
            \end{itemize}
            
            \item Satzbau {\small (sentence structure)}
            \begin{itemize}
                \item Wortstellung {\small (word order)}
                \begin{itemize}
                    \item Aussagesatz {\small (declarative sentence)}
                    \item Fragesatz {\small (interrogative sentence)}
                    \item Satzklammer {\small (sentence bracket)}
                    \item Stellung von Satzgliedern {\small (position of sentence elements)}
                \end{itemize}
                
                \item Satzarten {\small (types of sentences)}
                \begin{itemize}
                    \item Aussagesatz {\small (declarative sentence)}
                    \item Fragesatz {\small (interrogative sentence)}
                    \item Aufforderungssatz {\small (imperative sentence)}
                    \item Wunschsatz {\small (wish sentence)}
                    \item Ausrufesatz {\small (exclamatory sentence)}
                \end{itemize}
                
                \item Satzglieder {\small (sentence elements)}
                \begin{itemize}
                    \item Prädikat {\small (predicate)}
                    \item Subjekt {\small (subject)}
                    \item Akkusativobjekt {\small (accusative object)}
                    \item Dativobjekt {\small (dative object)}
                    \item Präpositionalobjekt {\small (prepositional object)}
                    \item Adverbiale {\small (adverbials)}
                \end{itemize}
                
                \item Negation {\small (negation)}
                \begin{itemize}
                    \item \textit{nicht} {\small (not)}
                    \item \textit{kein} {\small (no / none)}
                    \item pauschale Negation {\small (general negation)}
                    \item Negationswörter {\small (negative words)}
                \end{itemize}
            \end{itemize}
            
            \item Nebensätze und Satzverbindungen {\small (subordinate clauses and sentence connections)}
            \begin{itemize}
                \item Grundtypen {\small (basic types)}
                \begin{itemize}
                    \item dass-Sätze {\small (that-clauses)}
                    \item ob-Sätze {\small (if/whether-clauses)}
                    \item Fragesätze als Nebensätze {\small (indirect questions)}
                    \item Relativsätze {\small (relative clauses)}
                    \item Infinitivsätze mit \textit{zu} {\small (infinitive clauses with \textit{zu})}
                \end{itemize}
                
                \item adverbiale Nebensätze {\small (adverbial clauses)}
                \begin{itemize}
                    \item kausale Nebensätze {\small (causal clauses)}
                    \item konzessive Nebensätze {\small (concessive clauses)}
                    \item konditionale Nebensätze {\small (conditional clauses)}
                    \item modale Nebensätze {\small (modal clauses)}
                    \item temporale Nebensätze {\small (temporal clauses)}
                    \begin{itemize}
                        \item \textit{wenn} und \textit{als} {\small (when and when (past))}
                        \item \textit{seit[dem]} und \textit{bis} {\small (since and until)}
                        \item \textit{während} {\small (while)}
                        \item \textit{nachdem} {\small (after)}
                        \item \textit{bevor} und \textit{ehe} {\small (before)}
                        \item \textit{sobald} {\small (as soon as)}
                        \item \textit{solange} {\small (as long as)}
                    \end{itemize}
                    \item konsekutive Nebensätze {\small (consecutive/result clauses)}
                    \item adversative Nebensätze {\small (adversative clauses)}
                    \item Finalsätze {\small (purpose clauses)}
                    \begin{itemize}
                        \item finale Nebensätze {\small (finite purpose clauses)}
                        \item finale Infinitivsätze {\small (infinitive purpose clauses)}
                        \item \textit{damit}-Sätze {\small (so that clauses)}
                    \end{itemize}
                \end{itemize}
                
                \item Konjunktionen {\small (conjunctions)}
                \begin{itemize}
                    \item nebenordnende Konjunktionen {\small (coordinating conjunctions)}
                    \item unterordnende Konjunktionen {\small (subordinating conjunctions)}
                \end{itemize}
            \end{itemize}
            
            \item Passiv {\small (passive voice)}
            \begin{itemize}
                \item Vorgangspassiv {\small (process passive)}
                \begin{itemize}
                    \item mit Subjekt {\small (with subject)}
                    \item ohne Subjekt {\small (without subject)}
                    \item Zeiten im Vorgangspassiv {\small (tenses in process passive)}
                    \item Vorgangspassiv mit Modalverben {\small (process passive with modal verbs)}
                    \item Vorgangspassiv ohne Täter {\small (process passive without agent)}
                \end{itemize}
                
                \item Zustandspassiv {\small (state passive)}
            \end{itemize}
            
            \item Modalität {\small (modality)}
            \begin{itemize}
                \item objektive Bedeutungen der Modalverben {\small (objective meanings of modal verbs)}
                \begin{itemize}
                    \item Möglichkeit {\small (possibility)}
                    \item Fähigkeit {\small (ability)}
                    \item Erlaubnis {\small (permission)}
                    \item Verbot {\small (prohibition)}
                    \item Notwendigkeit {\small (necessity)}
                    \item Pflicht {\small (obligation)}
                    \item Wunsch {\small (wish)}
                    \item Wille {\small (will)}
                \end{itemize}
                
                \item subjektive Bedeutungen der Modalverben {\small (subjective meanings of modal verbs)}
                \begin{itemize}
                    \item Vermutung {\small (assumption)}
                    \item Unsicherheit {\small (uncertainty)}
                    \item Gerücht {\small (rumor)}
                    \item Empfehlung {\small (recommendation)}
                \end{itemize}
            \end{itemize}
            
            \item Umformungen und Stil {\small (transformations and style)}
            \begin{itemize}
                \item Nominalisierung {\small (nominalization)}
                \item Verbalisierung {\small (verbalization)}
                \item Umwandlung adverbialer Präpositionalgruppen in Nebensätze {\small (conversion of adverbial prepositional phrases into subordinate clauses)}
            \end{itemize}
        \end{itemize}
\end{document}